% !PRINTED 613
\weberpage{613}

Bilden wir nun nach (4) die Zusammensetzung
\[
\lambda \lambda_0^{-h} = \left(z,\, z + a_0^{-h} b + \frac{a_0^{-h} - 1}{a_0 - 1} b_0\right),
\]
so erhalten wir dann und nur dann die identische Substitution, wenn
\[
\frac{b_0}{a_0 - 1} \equiv \frac{b}{a_0^h - 1}\ (\mathrm{mod}\ n).
\]

Wenn diese Congruenz für \gesperrt{alle} Substitutionen $\lambda$ von $L$ erfüllt ist, so bleibt das Element $\dfrac{b_0}{a_0 - 1}$ durch die ganze Gruppe ungeändert, und $L$ ist intransitiv. Folglich giebt es in einer transitiven Gruppe $L$ immer ein Element $\lambda$, für das $\lambda \lambda_0^{-h}$ nicht die identische Substitution ist, obwohl der Index Null ist, und $L$ enthält also die ganze cyklische Gruppe.

Bildet man die verschiedenen Potenzen von $\lambda_0$, so ergiebt sich, dass der Exponent $h$ in den Substitutionen $\lambda$ der Gruppe $L$ jeden Werth annehmen kann. Ist $e$ der kleinste positive Exponent, der der Bedingung $a_0^e \equiv 1\ (\mathrm{mod}\ n)$ genügt, so kommen also die Werthe $h = 0, 1, 2 \ldots e - 1$ vor, und darin ist $e$ ein Theiler von $n - 1$. Aus der zusammengesetzten Substitution
\[
(z, a_0^h z + b)(z, z + 1) = (z, a_0^h z + b + 1),
\]
die ja auch in $L$ vorkommen muss, sieht man weiter, dass mit jedem Werth von $h$ jeder der Werthe $b = 0, 1, 2, \ldots n - 1$ verbunden vorkommt, und die ganze Gruppe $L$ besteht daher aus den Substitutionen
\weq{5}{\lambda = (z, a_0^h z + b), \quad \begin{array}{l} h = 0, 1 \ldots e - 1 \\ b = 0, 1 \ldots n - 1' \end{array}}
% [sic] source prints a raised apostrophe-like mark after "n-1" in the b-row
und ist vom Grade $en$. Umgekehrt bildet jedes System von Substitutionen dieser Form eine Gruppe.

Von den linearen Gruppen gilt nun der folgende Satz:

I. \gesperrt{Ist eine transitive lineare Gruppe $L$ Normaltheiler einer anderen Permutationsgruppe $P$ derselben $n$ Ziffern, so ist auch $P$ eine lineare Gruppe.}

Der Satz lässt sich so beweisen:
Wenn $\pi$ eine beliebige Permutation ist:
\[
\pi = \begin{pmatrix} 0, & 1, & 2 \ldots n - 1 \\ a_0, & a_1, & a_2 \ldots a_{n-1} \end{pmatrix},
\]
%CONT>>
